% Source: inputs/original_02.pdf p.1, 2013 New Curriculum Paper 1 liberal arts Q9 option C.
% Exact formula y=(1-\cos x)\sin x on [-pi,pi].
\usetikzlibrary{arrows.meta}
\begin{tikzpicture}[x=1.1cm,y=1.1cm,line width=.6pt,>={Stealth[length=2mm,width=1.4mm]},font=\normalsize]
\draw[->,line width=.6pt] (-3.5,0)--(3.79,0);
\draw[->,line width=.6pt] (0,-2.0)--(0,2.29);
\draw[line width=.6pt] (-3.1415926536,0)--(-3.1415926536,0.12);
\draw[line width=.6pt] (3.1415926536,0)--(3.1415926536,0.12);
\draw[line width=.6pt] (0,1)--(0.12,1);
\draw[domain=-3.1415926536:3.1415926536,samples=401,line width=.6pt] plot (\x,{(1-cos(\x r))*sin(\x r)});
\node[anchor=south east,inner sep=0pt,xshift=-1pt,yshift=2pt] at (0,2.29) {$y$};
\node[anchor=north west,inner sep=0pt,xshift=2pt,yshift=-2pt] at (3.79,0) {$x$};
\node[anchor=north west,inner sep=0pt,xshift=3pt,yshift=-4pt] at (0,0) {$O$};
\node[anchor=east,inner sep=0pt,xshift=-2pt,yshift=0pt] at (0,1) {$1$};
\node[anchor=south,inner sep=0pt,xshift=0pt,yshift=6pt] at (-3.142,0) {$-\pi$};
\node[anchor=north,inner sep=0pt,xshift=0pt,yshift=-5pt] at (3.142,0) {$\pi$};
\end{tikzpicture}
